% ------------------- %
% -- STANDARD CODE -- %
% ------------------- %

\begin{luacode*}
local nb_colors = #PALETTE

local final_file = io.open(
  "PROJECT-PALETTE.lua",
  "w"
)

local dev_file = io.open(
  "pal-def.lua",
  "r"
)

local in_header = false

local magic_comments = {}
local luadraw_code   = {}

for line in dev_file:lines() do
  if line:match("^%-%-%-%-%-%-+$") then
    in_header = not in_header

  elseif in_header then
    table.insert(magic_comments, line)

  else
    table.insert(luadraw_code, line)
  end
end

local luadraw_precode = [[--
--
-- Here is the luadraw code used.
--
-- lua::]]

local function est_vide(ligne)
  return ligne:match("^%s*$") ~= nil
end

while #luadraw_code > 0 and est_vide(luadraw_code[1]) do
  table.remove(luadraw_code, 1)
end

while #luadraw_code > 0 and est_vide(luadraw_code[#luadraw_code]) do
  table.remove(luadraw_code, #luadraw_code)
end

if #magic_comments > 0 then
  for i = 1, #magic_comments do
    final_file:write(magic_comments[i] .. "\n")
  end
end

final_file:write(luadraw_precode .. "\n")

for i = 1, #luadraw_code do
  final_file:write("--     " .. luadraw_code[i] .. "\n")
end

final_file:write("------\n\n")

dev_file:close()

final_file:write("PALETTE = {\n")

for n , rgb in ipairs(PALETTE) do
  final_file:write("  {")

  for i, c in ipairs(rgb) do
    final_file:write(c)

    if i ~= 3 then
      final_file:write(", ")
    end
  end

  final_file:write("}")

  if n ~= nb_colors then
    final_file:write(",")
  end

  final_file:write("\n")
end

final_file:write("}\n")

final_file:close()
\end{luacode*}


% ------------- %
% -- PALETTE -- %
% ------------- %

\centering

\bigskip

\begin{luadraw}{name = showcase-palette-dark}
local g = graph:new{
  window = {xmin, xmax, -5, 5},
  bbox   = false
}

g:Linewidth(10)

local A = Z(-20, 4)
local v = Z(0, -PALDIM)

for k = 1, PALSIZE do
  local color = rgb(PALETTE[k])

  g:Drectangle(
    A, A + PALDIM, A + PALDIM + v,
    "color = " .. COL_1 .. ", fill = " .. color
  )

  g:Dlabel(
    "\\Large$" .. k .. "$",
    A + Z(PALDIM, 1.15*PALDIM) / 2,
    {pos="S"}
  )

  A = A + PALDIM + PALDELTA
end

g:Show()
\end{luadraw}


% -------------- %
% -- SPECTRUM -- %
% -------------- %

\bigskip

\begin{luadraw}{name = showcase-spectrum-dark}
-- Let's work!
local g = graph:new{
  window = {xmin, xmax, -5, 5},
  bbox   = false
}

g:Linewidth(10)

local A  = Z(-5, 4)
local h  = Z(0, -PALDIM*1.30)
local dh = Z(0, -1.1)

local L = PALDIM * PALSIZE + PALDELTA * (PALSIZE - 1)
local N = 100

local dl = L / N

for k = 1, N do
  local color = palette(PALETTE, (k - 1) / (N - 1))

  g:Drectangle(
    A, A + h, A + h + dl,
    "color = " .. color .. ", fill = " .. color
  )

  A = A + dl
end

g:Drectangle(
  A, A + h, A + h - L,
  "color = " .. COL_1
)

g:Show()
\end{luadraw}


% ---------------------- %
% -- PRETTY NEFERTITI -- %
% ---------------------- %

\newpage

\begin{luadraw}{name = showcase-nefertiti-dark}
local obj_file = "./core/nefertiti.obj"

local polyhedron, bounding_box = read_obj_file(obj_file)

local g = graph3d:new{
  window3d = bounding_box,
  window   = {-5, 4, -7, 5},
  viewdir  = {35, 60},
  size     = {15, 15},
  bbox     = false
}

g:Scale(.6)

g:Dpoly(
  polyhedron,
  {
    usepalette = {PALETTE, "z" },
    mode       = mShadedOnly
  }
)

g:Show()
\end{luadraw}


% ------------------ %
% -- LEVEL SURFACE-- %
% ------------------ %

\vfill

\begin{luadraw}{name = showcase-level-surface-dark}
local cos, sin = math.cos, math.sin, math.pi

local g = graph3d:new{
  window3d = {0, 5, 0, 10, 0, 11},
  adjust2d = true,
  size     = {14.5, 9.5, 0},
  bbox     = false,
  viewdir  = {220, 60},
  margin   = {0, 0, 0, 0}
}

g:Scale(.65)

g:Linewidth(2)

local S = surface(
  function(u, v)
    return M(u, v, (u + v) / (2 + cos(u)*sin(v)))
  end,
  0, 5, 0, 10,
  {30, 30}
)

local n = 10

local colors = getpalette(PALETTE, n, true)

local niv, S1 = {}

for k = 1, n do
  S1, S = cutfacet(S, {M(0, 0, k), -vecK})

  insert(
    niv,
    {
      S1,
      {
        color     = colors[k],
        mode      = mShaded,
        edgewidth = 0.5
      }
    }
  )
end

insert(
  niv,
  {
    S,
    {color = colors[n + 1]}
  }
)

g:Linewidth(7)

g:Dboxaxes3d({
  grid      = true,
  gridcolor = COL_1,
  fillcolor = COL_2
})

g:Dmixfacet(table.unpack(niv))

for k = 1, n do
  g:Dballdots3d(
    M(5, 0, k),
    rgb(colors[k])
  )
end

g:Show()
\end{luadraw}


% ------------------------- %
% -- CONTINUOUS SPECTRUM -- %
% ------------------------- %

\newpage

\begin{luadraw}{name = showcase-contour-dark}
local g = graph:new{
  window = {-1, 6.5, -1.5, 11},
  size   = {16, 14, 0},
  margin = {0, 0, 0, 0},
  bbox   = false
}

g:Scale(.75)

local i = cpx.I

local sin, cos = math.sin, math.cos

local f = function(x, y)
  return (x + y) / (2 + cos(x)*sin(y))
end

local Lz = range(1, 10)

local colors = getpalette(PALETTE, 10)

g:Linewidth(3)

g:Dgradbox(
  {0, 5 + 10*i, 1, 1},
  {
    grid         = true,
    gridcolor    = COL_1,
    subgridcolor = COL_2
  }
)

g:Linewidth(20)

g:Dcontour(
  f,
  Lz,
  {
    view   = {0, 5, 0, 10},
    colors = colors
  }
)

for k = 1, 10 do
  local y = (2*k + 4) / 3*i

  g:Dseg(
    {5.25 + y, 5.5 + y},
    1,
    "color = " .. colors[k]
  )

  g:Labelcolor(colors[k])

  g:Dlabel(
    "$z = " .. k .. "$",
    5.5 + y,
    {pos = "E"}
  )
end

g:Show()
\end{luadraw}


% ---------------------------- %
% -- DELAUNAY TRIANGULATION -- %
% ---------------------------- %

\vfill

\begin{luadraw}{name = showcase-delaunay-dark}
local g = graph3d:new{
  bbox           = false,
  pictureoptions = "scale = 1.95"
}

g:Linewidth(6)

local i = cpx.I

local L = {
   0.285 + 1.46*i,
   1.556 - 0.142*i,
   2.344 + 1.313*i,
  -2.38  + 1.218*i,
   1.548 - 0.624*i,
   0.969 + 1.819*i,
  -0.086 - 2.191*i,
  -0.477 + 1.834*i,
  -0.904 + 1.322*i,
  -2.892 + 0.025*i
}

local T = delaunay(L)
local n = #T

local colors = getpalette(PALETTE,n)

for k = 1, n do
  g:Dpolyline(
    T[k],
    true,
       "color = " .. COL_2
    .. ", fill = " .. colors[k]
  )
end

g:Ddots(L, "color = " .. COL_2)

g:Show()
\end{luadraw}


% --------------- %
% -- 2D CHARTS -- %
% --------------- %

\vfill

\begin{luacode*}
function read_CSV(final_file, ignore_header)
  local data = {}

  for line in io.lines(final_file) do
    if ignore_header then
      ignore_header = false

    else
      local result = {}

      for part in string.gmatch(line, "([^,]+)") do
        table.insert(result, tonumber(part))
      end

      table.insert(data, result)
    end
  end

  return data
end
\end{luacode*}


\begin{luadraw}{name = showcase-chart-2D-dark}
local g = graph:new{
  window = {-5, 5, -4, 4},
}

g:Daxes(
  {0, 1, 1},
  {
    arrows       = "->",
    grid         = true,
    gridcolor    = COL_1,
    subgridcolor = COL_2
  }
)

local data = read_CSV(
  "./core/data-2D.csv",
  true
)

local i = cpx.I

local xyR

local colors = getpalette(PALETTE, 10)

for _, xyR in ipairs(data) do
  local col = colors[xyR[3]]

  g:Dcircle(
    xyR[1] + i*(xyR[2] -.5), .3 + .05*xyR[3],
    "fill = " .. col
  )
end

g:Show()
\end{luadraw}
